\documentclass[10pt,a4paper]{amsart}
\usepackage[margin=19mm]{geometry}
\usepackage{amsmath,amssymb,mathtools}
\usepackage[scr=rsfs,cal=euler]{mathalfa}
\usepackage{xfrac}
\usepackage[unicode,hidelinks,bookmarks=false]{hyperref}
\newcommand{\ie}{i.e.\ }
\newcommand{\cf}{cf.\ }
\newcommand{\resp}{respectively\ }
\newcommand{\Subsection}{\subsection}
\providecommand{\nobreakdash}{\nobreak\hbox{-}}
\setlength{\emergencystretch}{2em}
\multlinegap0pt
\usepackage{amssymb}
\usepackage{mathtools}
\usepackage{xcolor}
\usepackage[shortlabels]{enumitem}
\newtheorem{theorem}{Theorem}
\newcommand{\F}{\mathbb{F}}
\newcommand{\rank}{\mathrm{rank}}
\title{Formalizing building-up constructions of self-dual codes through isotropic lines in Lean}
\author{Jae-Hyun Baek, Jon-Lark Kim}
\date{Source: arXiv:2604.08485v2 (CC BY 4.0)}
\begin{document}
\maketitle

\noindent\emph{Source and licence.} Formalizing building-up constructions of self-dual codes through isotropic lines in Lean. Version arXiv:2604.08485v2 (CC BY 4.0), distributed by arXiv under the Creative Commons Attribution 4.0 International licence, http://creativecommons.org/licenses/by/4.0/, as recorded in the arXiv licence metadata for this exact version and shown on the abstract page https://arxiv.org/abs/2604.08485v2. Full licence notice: \texttt{licenses/arxiv-2604.08485-CC-BY-4.0.txt}.\par\smallskip

\noindent\textit{Editorial scope.} Counted unit: the article's theorem thm:qary-universal-rank-r (source0.tex lines 1448--1491). The article's complete original proof of it is delivered with it (source0.tex lines 1493--1605, one \texttt{proof} environment of 113 lines). No mathematics is written, repaired or re-derived: the statement and the proof are the article's own lines, in the article's own notation, and no retained line is substituted or re-flowed.  No further source-local context is needed: every cross-reference of every retained line resolves inside the two delivered spans.  Nothing was omitted from any retained span, so the package carries no omission record.  No retained line contains a citation, so the wrapper carries no bibliography and no bibliography entry.  No retained line contains a figure environment, an image-inclusion command or drawing code, so no image is delivered and the package declares no asset dependency.  Editorial formatting only, in addition: an AMS \texttt{amsart} preamble that restates the article's own declarations and macros used by the retained lines, namely the environments \texttt{theorem} on separate counters, together with the standard packages those lines need.  The retained environments print in this document as Theorem 1 in order of appearance, whereas the article prints the same items with the numbering of its own full document.  The complete native source of the article as distributed in the arXiv e-print for this exact version is bundled unchanged as \texttt{sources/198158/source0.tex}.
\begin{theorem}[Universal rank-\(r\) boxed normal form over
\texorpdfstring{$\F_q$}{Fq}]
\label{thm:qary-universal-rank-r}
Let \(K=\F_q\), where \(q\equiv1\pmod4\), fix \(c\in K^\times\) with
\(c^2=-1\), and fix an ordered pairing \(K^{2n}=(K^2)^n\).  For every
Euclidean self-dual code \(\mathcal C\subseteq K^{2n}\), put
\[
 U_c=\bigoplus_{j=1}^n K(1,c),\qquad
 r=\dim(\mathcal C\cap U_c),\qquad k=n-r.
\]
After a permutation of the \(n\) two-coordinate blocks and elementary row
operations, \(\mathcal C\) has a generator matrix
\begin{equation}
G_r(c;Q,\ell,D)=
\left[
\begin{array}{c|c}
\bigl(Q_{ij}(1,c)+\delta_{ij}(0,1)\bigr)_{k\times k}
  & \bigl(\ell_{it}\bigr)_{k\times r}\\ \hline
\bigl(-(D\partial_c\ell_j)_s(1,c)\bigr)_{r\times k}
  & \bigl(D_{st}(1,c)\bigr)_{r\times r}
\end{array}
\right],
\label{eq:universal-rank-r-box}
\end{equation}
where \(1\le i,j\le k\), \(1\le s,t\le r\),
\(\ell_{it}=(u_{it},v_{it})\in K^2\),
\[
 (\partial_c\ell_i)_t=v_{it}-cu_{it},\qquad
 Q_{ii}=\frac c2\bigl(1+\ell_i\cdot\ell_i\bigr),
\]
\(D\in M_r(K)\) is nonsingular, and the remaining coefficients satisfy
\begin{equation}
 c(Q_{ij}+Q_{ji})+\ell_i\cdot\ell_j=0
\qquad(i\ne j).
\label{eq:universal-rank-r-relation}
\end{equation}
The last \(r\) rows span \(\mathcal C\cap U_c\).  Conversely, every matrix
\eqref{eq:universal-rank-r-box} with \(\det D\ne0\) and
\eqref{eq:universal-rank-r-relation} has \(k+r\) independent, pairwise
orthogonal rows and generates a Euclidean self-dual code.  Moreover,
retaining any subset of the first \(k\) rows together with the corresponding
block columns and all terminal rows and columns gives a self-dual code of
the same form with the same \(D\).
\end{theorem}

\begin{proof}
For \(v=(v_1,\ldots,v_n)\in(K^2)^n\), write each block uniquely as
\[
 v_j=(\alpha_j(v),c\alpha_j(v)+\beta_j(v)),\qquad
 \beta_j(v)=v_{j,2}-cv_{j,1},
\]
and define \(\partial:\mathcal C\to K^n\) by
\(\partial(v)=(\beta_1(v),\ldots,\beta_n(v))\).  Since
\(\ker\partial=\mathcal C\cap U_c\), rank--nullity gives
\(\rank\partial=k\).  Choose \(k\) coordinates on which
\(\operatorname{im}\partial\) projects isomorphically, permute them to the
first \(k\) positions, and choose rows \(g_1,\ldots,g_k\in\mathcal C\) whose defects
restrict there to the standard basis.  Complete them by a basis
\(S_1,\ldots,S_r\) of \(\ker\partial\).  Applying \(\partial\) to a linear
relation first determines the coefficients of the \(g_i\); independence of
the \(S_s\) then shows that these \(n\) rows form a basis of \(\mathcal C\).

Let \(Q_{ij}\) be the first coordinate of the \(j\)-th pivot block of
\(g_i\), and write its terminal block-row as \(\ell_i\).  Thus its
\(j\)-th pivot block is
\(Q_{ij}(1,c)+\delta_{ij}(0,1)\).  Since every \(S_s\) lies in
\(\ker\partial\), write its pivot and terminal blocks as
\(A_{sj}(1,c)\) and \(D_{st}(1,c)\), respectively.  The identity
\[
 (\alpha,c\alpha+\beta)\cdot(\gamma,c\gamma+\delta)
   =c(\alpha\delta+\beta\gamma)+\beta\delta
\]
and the orthogonality of the chosen basis give
\[
 A_{si}+(D\partial_c\ell_i)_s=0,
\qquad
 \delta_{ij}+c(Q_{ij}+Q_{ji})+\ell_i\cdot\ell_j=0.
\]
The first equation gives the lower-left block in
\eqref{eq:universal-rank-r-box}; the second gives
\eqref{eq:universal-rank-r-relation}, and its diagonal part gives the stated
formula for \(Q_{ii}\).  It remains to prove that \(D\) is nonsingular.
Let \(\lambda=(\lambda_1,\ldots,\lambda_r)\in K^r\) satisfy
\(\lambda D=0\).  For every \(i\), the first identity yields
\[
 (\lambda A)_i
   =\sum_{s=1}^r\lambda_sA_{si}
   =-\sum_{t=1}^r(\partial_c\ell_i)_t
                    \sum_{s=1}^r\lambda_sD_{st}=0.
\]
Consequently, every pivot and terminal block of
\(\sum_s\lambda_sS_s\) is zero.  Since the \(S_s\) are linearly
independent, \(\lambda=0\).  Thus \(\ker D=0\), and hence
\(\det D\ne0\).

Conversely, suppose that \(D\) is nonsingular and that
\eqref{eq:universal-rank-r-relation} holds.  Denote the first \(k\) rows of
\eqref{eq:universal-rank-r-box} by \(p_1,\ldots,p_k\), put
\(q_{it}=(\partial_c\ell_i)_t\), and retain the notation
\(S_1,\ldots,S_r\) for the terminal rows.  Since
\((1,c)\cdot(1,c)=0\), one has
\[
 S_s\cdot S_u=0,
 \qquad
 p_i\cdot S_s
   =c\left(A_{si}+\sum_{t=1}^r q_{it}D_{st}\right)=0.
\]
Moreover,
\[
 p_i\cdot p_j=
 \begin{cases}
  1+2cQ_{ii}+\ell_i\cdot\ell_i=0,&i=j,\\
  c(Q_{ij}+Q_{ji})+\ell_i\cdot\ell_j=0,&i\ne j.
 \end{cases}
\]
Thus the displayed rows are pairwise orthogonal.

Suppose now that
\[
 \sum_{i=1}^k\xi_i p_i+\sum_{s=1}^r\mu_sS_s=0.
\]
Applying \((x,y)\mapsto y-cx\) to the \(j\)-th pivot block gives
\(\xi_j=0\), because the defect of that block is \(\delta_{ij}\) in
\(p_i\) and zero in every \(S_s\).  Hence \(\xi_1=\cdots=\xi_k=0\).
The first coordinates of the terminal blocks then give
\[
 0=\sum_{s=1}^r\mu_sD_{st}=(\mu D)_t
 \qquad(1\le t\le r).
\]
Since \(D\) is nonsingular, \(\mu=0\).  The \(k+r\) rows are therefore
linearly independent and pairwise orthogonal in \(K^{2(k+r)}\); their row
space is consequently self-dual.  If
\(x=\sum_i\xi_ip_i+\sum_s\mu_sS_s\) belongs to \(U_c\), then the defect of
its \(j\)-th pivot block is \(\xi_j\), so \(\xi_j=0\) for every \(j\).
Conversely, every linear combination of the \(S_s\) lies in \(U_c\).
Therefore, writing
\(\mathcal C_r:=\langle G_r(c;Q,\ell,D)\rangle_K\),
\[
 \mathcal C_r\cap U_c
   =\langle S_1,\ldots,S_r\rangle_K.
\]

Finally, let \(I\subseteq\{1,\ldots,k\}\).  Retaining the rows \(p_i\)
with \(i\in I\), their pivot block columns, and all terminal rows and block
columns replaces \(Q\) by \(Q|_{I\times I}\) and \((\ell_i)_{i=1}^k\) by
\((\ell_i)_{i\in I}\), while leaving \(D\) unchanged.  The two Gram
identities above restrict to indices in \(I\).  Hence its rows are pairwise
orthogonal; the pivot-block defects determine the coefficients of the
retained \(p_i\), and the nonsingularity of \(D\) determines those of the
\(S_s\).  Thus the restricted matrix has \(|I|+r\) independent rows in
\(K^{2(|I|+r)}\) and generates a self-dual code.  For
\(I=\varnothing\), this row space is
\[
 \{(z_1(1,c),\ldots,z_r(1,c)):z\in K^r\},
\]
the full terminal isotropic space in \(K^{2r}\), because \(D\) is
nonsingular.  This proves the final assertion.
\end{proof}

\end{document}
